\chapter{模板使用指南}\label{chap:guide}

本文档是 \TongjiThesis{} 模板的使用示例，基本覆盖了模板的全部排版功能。建议在阅读本示例的同时参考源代码，了解各排版效果的实现方式。在正式撰写毕业设计（论文）时，请删除本章及后续示例章节，替换为实际内容。

后续示例章节依次演示：\cref{chap:introduction}（标题、列表、字体）、\cref{chap:float}（图、表、算法、代码）、\cref{chap:math}（数字单位、公式、定理）、\cref{chap:reference}（参考文献、脚注、交叉引用）、\cref{chap:conclusion}（总结与展望），附录则示范 \cs{appendixsection}/\cs{appendixsubsection} 的用法。

\section{文档类选项}

在 \verb|main.tex| 的 \verb|\documentclass| 中设置选项。本模板基于 \pkg{ctexbook}，未被识别的选项会自动传递给 \pkg{ctexbook}。

\begin{table}[htbp]
  \centering
  \caption{文档类选项一览}\label{tab:class-options}
  \begin{tabularx}{\linewidth}{l l X}
    \toprule
    \textbf{选项} & \textbf{默认值} & \textbf{说明} \\
    \midrule
    \texttt{oneside / twoside} & \texttt{oneside} & 单面或双面排版。双面模式下章节从奇数页开始，装订线自动左右交替。 \\
    \texttt{degree} & \texttt{bachelor} & 学位类型，目前支持 \texttt{bachelor}（学士）。 \\
    \texttt{field} & \texttt{science} & 专业类别。\texttt{science}：理工科（工科/理科），章节编号采用阿拉伯数字体系；\texttt{humanities}：文科（经管/人文/法学/外语/艺术），章节编号采用汉字体系，详见\cref{tab:heading-levels}。 \\
    \texttt{fontset} & \texttt{fandol} & 字体集。\texttt{fandol}（跨平台，默认）、\texttt{windows}、\texttt{mac}、\texttt{founder}、\texttt{adobe}。详见 \verb|style/font/| 目录下的字体定义文件。 \\
    \texttt{times} & \texttt{false} & \texttt{true}：通过 \pkg{fontspec} 调用系统 Times New Roman 字体；\texttt{false}：使用 \pkg{newtx} 宏包提供的 \TeX{} 原生 Times 风格字体。两者视觉效果相近。 \\
    \texttt{fullwidthstop} & \texttt{circle} & 句号样式：\texttt{circle}（默认，保留“。”）或 \texttt{dot}（替换为“．”全角句点）。 \\
    \texttt{minted} & \texttt{false} & 代码高亮方案。\texttt{false} 使用 \pkg{listings}（默认，无外部依赖），\texttt{true} 使用 \pkg{minted}（需安装 Python 3.11–3.13 和 Pygments）。 \\
    \texttt{biblatex} & \texttt{true} & 参考文献方案。\texttt{true} 使用 \BibLaTeX{} + \texttt{biber}（功能更丰富），\texttt{false} 使用 \BibTeX{} + \pkg{gbt7714}（编译更快）。 \\
    \texttt{algo} & \texttt{algpseudocode} & 算法排版方案。\texttt{algpseudocode}（默认）与 \texttt{algorithm2e} 语法不同，详见说明。 \\
    \texttt{doctype} & \texttt{thesis} & 文档类型。\texttt{thesis}（默认）：本指南其余部分描述的论文正文；\texttt{taskbook}/\texttt{proposal}/\texttt{midterm}：分别生成任务书/开题报告/中期报告，详见\cref{sec:doctypes}。 \\
    \bottomrule
  \end{tabularx}
\end{table}

\section{填写论文信息}

论文的封面、信息说明页与中英文摘要的个人信息分散在 \verb|chapters/| 目录下的两个文件中。各文件均已提供带注释的示例，照搬并替换字段值即可；模板据此自动生成对应页面。

\begin{itemize}
  \item \textbf{封面与信息说明页} —— \verb|chapters/metadata.tex|：\cs{school}、\cs{major}、\cs{student}、\cs{thesistitle}、\cs{thesistitleeng}、\cs{thesisadvisor}、\cs{thesisdate}（封面），以及 \cs{infotype}（\texttt{thesis}、\texttt{design} 或 \texttt{engineering}）、\cs{infoabstract}、\cs{infothesiswords}（毕业论文）或 \cs{infodrawings} + \cs{infowordcount}（毕业设计）、\cs{infomaterials}（信息说明页）。封面与信息说明页由随后的 \cs{MakeCover}、\cs{MakeInfoPage} 命令自动生成。
  \item \textbf{中英文摘要} —— \verb|chapters/00_abstract.tex|：\cs{MakeAbstract}\verb|{正文}{关键词}| 与 \cs{MakeAbstractEng}\verb|{Text}{Keywords}|，以及可选的摘要页论文标题 \cs{abstracttitle} 与 \cs{abstracttitleeng}。
\end{itemize}

可在 \cs{thesistitle} 与 \cs{thesistitleeng} 中使用 \verb|\\| 为封面标题手动添加换行，信息说明页标题会自动移除这些换行。若中英文摘要页需要使用独立换行位置，可在 \verb|chapters/metadata.tex| 中设置 \cs{abstracttitle}\verb|{摘要中文标题}{摘要中文副标题}| 与 \cs{abstracttitleeng}\verb|{Abstract English Title}{Abstract English Subtitle}|；未设置 \cs{abstracttitle} 与 \cs{abstracttitleeng} 时，中英文摘要页论文标题会沿用 \cs{thesistitle} 与 \cs{thesistitleeng}，并自动移除这些换行。

\section{文档结构}

\verb|main.tex| 已预置了完整的论文结构，用户只需修改各 \verb|\input| 文件的内容：

\begin{table}[htbp]
  \centering
  \caption{文档结构一览}\label{tab:doc-structure}
  \begin{tabularx}{\linewidth}{l l X}
    \toprule
    \textbf{结构} & \textbf{命令} & \textbf{说明} \\
    \midrule
    封面       & \verb|\MakeCover|      & 由 \verb|metadata.tex| 中的信息自动生成 \\
    信息说明页 & \verb|\MakeInfoPage|   & 由 \verb|metadata.tex| 中的信息自动生成 \\
    前置部分   & \verb|\frontmatter|    & 页码切换为大写罗马数字（I, II, III\ldots） \\
    摘要     & \verb|\MakeAbstract|     & 中英文摘要及关键词（\verb|00_abstract.tex|） \\
    目录     & \verb|\tableofcontents|  & 自动生成；可取消注释 \verb|\listoffigures|、\verb|\listoftables| \\
    正文     & \verb|\mainmatter|       & 页码切换为阿拉伯数字，从1重新开始 \\
    参考文献 & \verb|\makereferences|   & 自动排版，格式符合 GB/T 7714 \\
    附录     & \verb|\appendix|         & 新增「附录」章页，节编号为 A、B、C（\verb|appendix.tex|） \\
    谢辞     & \verb|\backmatter|       & 取消章编号，页码继续（\verb|ack.tex|） \\
    \bottomrule
  \end{tabularx}
\end{table}

正文各章放在 \verb|chapters/| 下，通过 \verb|main.tex| 中的 \verb|\input{chapters/XX_name}| 注册；参考文献数据库文件在导言区通过 \verb|\tjbibresource{bib/note.bib}| 指定，支持逗号分隔的多个文件。

\section{任务书 / 开题报告 / 中期报告}\label{sec:doctypes}

除本文档（论文正文，\texttt{doctype=thesis}，默认）外，模板还提供 3 个独立的官方行政文档，通过 \texttt{doctype} 选项选择：

\begin{table}[htbp]
  \centering
  \caption{doctype 选项一览}\label{tab:doctype-options}
  \begin{tabularx}{\linewidth}{l l l X}
    \toprule
    \textbf{doctype} & \textbf{文档} & \textbf{入口文件} & \textbf{说明} \\
    \midrule
    \texttt{thesis}   & 论文正文（默认） & \verb|main.tex|     & 完整封面/信息页/前置/正文/参考文献/附录/谢辞结构，详见\cref{tab:doc-structure}。 \\
    \texttt{taskbook} & 任务书           & \verb|taskbook.tex| & 课题信息表 + 起讫时间（由 \cs{taskbookperiod} 自动换算周数）+ 课题背景/主要内容/基本要求及成果形式/进度安排/资料及参考文献/其他要求 + 签字区。 \\
    \texttt{proposal} & 开题报告         & \verb|proposal.tex| & 课题信息表 + 填表日期 + 课题背景/方案介绍/参考文献/审核意见。 \\
    \texttt{midterm}  & 中期报告         & \verb|midterm.tex|  & 课题信息表 + 填表日期 + 填写说明 + 进展情况/成果/问题/计划/审核意见。 \\
    \bottomrule
  \end{tabularx}
\end{table}

三者与正文共用 \verb|chapters/metadata.tex| 中的 \cs{school}、\cs{major}、\cs{student}、\cs{thesistitle}、\cs{thesisadvisor}，无需重复填写。各入口文件结构一致，均为 \cs{MakeDocument}（依据 \texttt{doctype} 分派到 \cs{MakeTaskBook}、\cs{MakeProposal} 或 \cs{MakeMidterm}，仅渲染封面表格与签字区，如 \cs{MakeCover} 之于论文封面）加上对应的正文内容文件：

\begin{itemize}
  \item \texttt{taskbook.tex} \cs{MakeDocument} 后 \verb|\tjformsection{一、\tjheadertext 的课题背景}
\tjformnote{根据实际情况写为毕业设计或毕业论文，下同（此项成稿删除）}
\tjformhint{请在此处填写课题的背景与来源，说明选题的实际意义或理论意义。}

\tjformsection{二、\tjheadertext 应完成的主要内容}
\tjformhint{请列出毕业设计（论文）应完成的主要研究内容或设计内容。}

\tjformsection{三、\tjheadertext 的基本要求及应完成的成果形式}
\tjformhint{请说明基本要求，以及应完成的成果形式（如论文、图纸、软件、模型等）。}

\tjformsection{四、\tjheadertext 的进度安排}
\tjformhint{请列出各阶段的起止时间与主要工作内容，可使用列表或表格呈现。}

\tjformsection{五、\tjheadertext 应收集的资料及主要参考文献}
\tjformhint{请列出需要收集的资料及主要参考文献。}

\tjformsection{六、其他要求}
\tjformnote{（此项根据实际情况为可选项）}

\vfill
% 任务书正文末尾的签字行为四号黑体（审核意见框内的签字行为小四号，见开题/中期报告）
\tjformsignature[\zihao{4}]{指导教师签名}{9em}
|
  \item \texttt{proposal.tex} \cs{MakeDocument} 后 \verb|\tjformsection{一、\tjheadertext 课题背景}
\tjformhint{2 级标题：小四号，黑体，顶格，序号与题名之间空一格，正文\ 宋体小4号字，行距1.5倍，段前0.5行，段后0.5行。首行缩进2字符，参考文献悬挂0.74cm。}
\tjformnote{根据实际情况写为毕业设计或毕业论文，下同（此项成稿删除）}

\begin{enumerate}[label=\arabic*．,leftmargin=2em,labelsep=0pt]
  \item 课题来源及研究的目的和意义

  请在此处填写课题来源及研究的目的和意义。

  \item 国内外在该方向的研究现状和发展趋势

  请在此处填写国内外研究现状和发展趋势。
\end{enumerate}

\tjformsection{二、\tjheadertext 方案介绍}

\begin{enumerate}[label=\arabic*．,leftmargin=2em,labelsep=0pt]
  \item 主要研究内容（包括但不限于预期研究目标、研究内容要点等。）

  请在此处填写主要研究内容。

  \item 研究方案（包括但不限于拟采取的研究方法、技术路线图、可行性分析等）

  请在此处填写研究方案。

  \item 工作进度安排

  请在此处填写工作进度安排。
\end{enumerate}

\tjformsection{三、\tjheadertext 的主要参考文献}

请在此处列出主要参考文献（悬挂缩进 0.74cm，与正文参考文献格式一致）。

\tjformsection{四、审核意见}

\begin{tjreviewbox}
  指导教师意见：\tjformnote{（课题难度是否适中、工作量是否饱满、进度安排是否合理、工作条件是否具备、是否同意开题等）}

  \vspace{5\baselineskip}
  \begin{flushright}
    指导教师签名：\ulinebox{8em}\qquad\tjdateblank
  \end{flushright}
\end{tjreviewbox}

\begin{tjreviewbox}
  专业审核意见：\tjformnote{（适当对各项内容给与意见和评价，字数不限。）}

  \vspace{3\baselineskip}
  \begin{flushright}
    负责人签名：\ulinebox{8em}\qquad\tjdateblank
  \end{flushright}
\end{tjreviewbox}
|
  \item \texttt{midterm.tex} \cs{MakeDocument} 后 \verb|\tjformnotice{填写说明}

\begin{tjformnoticebody}
\begin{enumerate}[label=\arabic*.,leftmargin=2em,labelsep=0.5em]
  \item 请每位学生根据学校及学院检查的要求认真进行自查，及时发现课题研究过程中存在的问题，分析原因，并提出解决思路和措施，明确下一阶段任务。
  \item 每位学生应根据项目实施情况认真、实事求是填写。
  \item 该表填写完毕后，须请指导教师审核，并签署意见。
  \item 《同济大学本科生毕业设计（论文）中期检查报告》将作为答辩资格审查的主要材料之一。
\end{enumerate}
\end{tjformnoticebody}

\clearpage

\tjformsection{一、课题进展情况}
\tjformhint{是否按开题报告的内容及进度安排进行}
\tjformhint{请在此处填写课题进展情况。}

\tjformsection{二、课题研究已取得的阶段性成果：}
\tjformhint{已完成的研究工作及成果}
\tjformhint{请在此处填写已取得的阶段性成果。}

\tjformsection{三、存在的问题与解决思路}
\tjformhint{请在此处填写存在的问题与解决思路。}

\tjformsection{四、下一阶段的工作计划和研究内容}
\tjformhint{请在此处填写下一阶段的工作计划和研究内容。}

\tjformsection{五、审核意见}

\begin{tjreviewbox}
  \tjformlabel{指导教师意见：}\tjformhint{（建议对各项内容给与意见和评价，字数不限。）}

  \vspace{7\baselineskip}
  \tjformsignature{指导教师签名}{8em}
\end{tjreviewbox}

\begin{tjreviewbox}
  \tjformlabel{专业审核意见：}

  \vspace{1.5\baselineskip}
  % 默认两格均留空，供打印后手写勾选；如需提前在电子版中标注，在
  % chapters/metadata.tex 中调用 \midtermdecision{agree} 或 \midtermdecision{disagree}
  \tjcheckboxeq{\tongjimidtermdecision}{agree}~同\ 意\qquad\tjcheckboxeq{\tongjimidtermdecision}{disagree}~不\ 同\ 意

  \vspace{5\baselineskip}
  \tjformsignature{负责人签名}{8em}
\end{tjreviewbox}
|
\end{itemize}

正文分节内容（课题背景、进度安排等）由上述 \verb|chapters/*_body.tex| 提供，可直接编辑替换示例文字，用法与 \verb|appendix.tex| 类似。任务书的起讫周数通过 \verb|\taskbookperiod{起始年}{起始月}{起始日}{截止年}{截止月}{截止日}|（在 \verb|chapters/metadata.tex| 中设置）按天数自动换算（向上取整），无需手动填写周数。

构建命令：\verb|make taskbook| / \verb|make proposal| / \verb|make midterm|（或一次性 \verb|make forms|），Windows 下对应 \verb|.\make.bat taskbook| 等。

\section{常用命令速查}

\tabref{tab:template-commands}列出了本模板提供的常用命令，具体用法在后续各章中演示。

\begin{table}[htbp]
  \centering
  \caption{模板常用命令}\label{tab:template-commands}
  \begin{tabularx}{\linewidth}{l X}
    \toprule
    \textbf{命令} & \textbf{用途} \\
    \midrule
    \multicolumn{2}{l}{\textsf{交叉引用}} \\
    \verb|\cref{label}|           & 智能引用，自动生成“第~X~章”“图~X”“式~(X)”等 \\
    \verb|\chapref{label}|        & 引用章：第~X~章 \\
    \verb|\secref{label}|         & 引用节：第~X~节 \\
    \verb|\figref{label}|         & 引用图：图~X \\
    \verb|\tabref{label}|         & 引用表：表~X \\
    \verb|\eqref{label}|          & 引用公式：式~(X) \\
    \midrule
    \multicolumn{2}{l}{\textsf{参考文献}} \\
    \verb|\tjbibresource{file}|   & 导言区指定 \texttt{.bib} 文件 \\
    \verb|\makereferences|        & 输出参考文献列表 \\
    \verb|\cite{key}|             & 上标引用 \\
    \midrule
    \multicolumn{2}{l}{\textsf{排版辅助}} \\
    \verb|\pkg{name}|             & 排版宏包名称（带 CTAN 链接） \\
    \verb|\cs{name}|              & 排版命令名称（如 \cs{section}） \\
    \verb|\TJCircled{n}|           & 带圈数字：\TJCircled{1} \TJCircled{2} \TJCircled{3}（0–50 Unicode，51+ TikZ 画圈） \\
    \verb|\TJCircledLetter{x}|     & 带圈字母：\TJCircledLetter{A} \TJCircledLetter{b}（A–Z / a–z，或计数器名自动检测） \\
    \verb|\TongjiThesis|            & 项目 Logo：\TongjiThesis \\
    \verb|\dd|, \verb|\ee|, \verb|\ii|, \verb|\jj| & 直立体微分 $\dd$、自然对数底 $\ee$、虚数单位 $\ii$、$\jj$ \\
    \bottomrule
  \end{tabularx}
\end{table}
